\documentclass[11pt]{article}
\usepackage{mathunal-base}
\mathunalfuente{serif}

\renewcommand{\examtitulo}{Supletorio Tercer Examen Parcial de Métodos Numéricos}
\renewcommand{\examfecha}{Diciembre 3 de 2007}

\begin{document}

\examheader

\vspace{4pt}
\noindent\textbf{Duración: 1 hora y 50 minutos}

\vspace{4pt}
\noindent Nombre: \dotfill\ Carné: \dotfill\ Grupo: \dotfill

\begin{center}\textbf{Toda respuesta debe estar debidamente justificada}\end{center}

\begin{enumerate}
\item{} Considere el P.V.I.
\[
\left\{
\begin{aligned}
y' &= \frac{4t^2y}{1+t^4}, \qquad 0\leq t\leq 2\\
y(0) &= 1
\end{aligned}
\right.
\]
\begin{enumerate}
\item{} (15\%) Demuestre que la función $f(t,y)=\dfrac{4t^2y}{1+t^4}$ satisface una condición de Lipschitz para $0\leq t\leq 2$. Determine el menor valor de la constante de Lipschitz.

\vspace{4cm}

\item{} (5\%) ¿Tiene este P.V.I. una única solución?

\vspace{2cm}
\end{enumerate}

\item{} Considere el P.V.I.
\[
\left\{
\begin{aligned}
t^2y''-2ty'+2y &= t^3\ln(t)\\
y(1)&=1,\quad y'(1)=0
\end{aligned}
\right.
\]
\begin{enumerate}
\item{} (10\%) Convierta este P.V.I. a la forma vectorial
\[
\left\{
\begin{aligned}
U' &= F(t,U)\\
U(t_0) &= U_0
\end{aligned}
\right.
\]
\noindent con $U$ y $U_0$ vectores de $\mathbb{R}^2$.

\vspace{4cm}

\item{} (20\%) Utilice el método de Taylor de orden 2 para aproximar $y(1.5)$ utilizando tamaño de paso $h=0.5$.

\vspace{5cm}
\end{enumerate}

\item{} Considere el problema de contorno o de valores en la frontera
\[
\left\{
\begin{aligned}
y''-\frac{y'}{t}-\frac{3y}{t^2}-\frac{\ln(t)}{t} &= -1, \qquad 1\leq t\leq 2\\
y(1) &= 0,\quad y(2)=0
\end{aligned}
\right.
\]
\begin{enumerate}
\item{} (8\%) ¿Tiene este problema de contorno o de valores en la frontera una única solución?

\vspace{3cm}

\item{} (12\%) Las soluciones aproximadas de los problemas de valor inicial asociados, ambos para $1\leq t\leq 2$,
\[
\left\{
\begin{aligned}
w''-\frac{w'}{t}-\frac{3w}{t^2}-\frac{\ln(t)}{t} &= -1\\
w(1)&=0,\ w'(1)=0
\end{aligned}
\right.
\qquad\qquad
\left\{
\begin{aligned}
u''-\frac{u'}{t}-\frac{3u}{t^2} &= 0\\
u(1)&=0,\ u'(1)=1
\end{aligned}
\right.
\]
\noindent están dadas en la siguiente tabla

\vspace{6pt}
\begin{center}
\begin{tabular}{|c|c|c|c|c|c|}
\hline
 & $t=1$ & $t=1.2$ & $t=1.4$ & $t=1.6$ & $t=1.8$ \\ \hline
$w(t)$ & $0$ & $-0.0202$ & $-0.0831$ & $-0.1938$ & $-0.3590$ \\ \hline
$u(t)$ & $0$ & $0.2236$ & $0.5073$ & $0.8676$ & $1.3188$ \\ \hline
\end{tabular}
\end{center}
\noindent (con $w(2)=-0.5861$ y $u(2)=1.8746$.)

\noindent Utilice esta tabla para aproximar la solución del problema de contorno por el método del disparo lineal con tamaño de paso $h=0.2$.

\vspace{6cm}
\end{enumerate}

\item{} Considere el problema
\[
\left\{
\begin{aligned}
\frac{\partial u}{\partial t}-\frac{1}{4}\frac{\partial^2 u}{\partial x^2} &= 0, \qquad 0<x<2,\ t>0\\
u(x,0) &= \operatorname{sen}\left(\frac{\pi}{4}x\right)-\cos\left(\frac{\pi}{2}x\right)+1, \qquad 0\leq x\leq 2\\
u(0,t) &= t, \qquad t\geq 0\\
u(2,t) &= 2+e^{-t}, \qquad t\geq 0
\end{aligned}
\right.
\]
\begin{enumerate}
\item{} (6\%) ¿Es el método de diferencia progresiva con $h=0.5$ y $k=0.2$, estable?

\vspace{3cm}

\item{} (24\%) Hallar una aproximación a la solución $u(x,t)$ para $t=0.4$ empleando el método de diferencia progresiva con $h=0.5$ y $k=0.2$.

\vspace{6cm}
\end{enumerate}
\end{enumerate}

\examfooter
\end{document}
